% !PRINTED 134
\weberpage{134}
%<<CONT
\weq{5}{
f_2\!\left(\frac{\gamma+\delta\omega}{\alpha+\beta\omega}\right) = \left(\frac{2}{\beta}\right)\varrho\, e^{\frac{3\pi i}{8}\beta(2\delta-\alpha)}\, f_1(\omega), \qquad \alpha \equiv 0 \pmod{2},
}
and if in this we replace $\omega$ by $\omega+3$ and $\gamma, \alpha$ by $\gamma-3\delta, \alpha-3\beta$, it follows that
\weq{6}{
f_2\!\left(\frac{\gamma+\delta\omega}{\alpha+\beta\omega}\right) = -\varrho\, e^{\frac{3\pi i}{8}\beta(2\delta-\alpha)}\, f(\omega), \qquad \alpha-\beta \equiv 0 \pmod{2}.
}

If in (4), (5), (6) we set
\[
f_2\!\left(\frac{\gamma+\delta\omega}{\alpha+\beta\omega}\right) = f_1\!\left(-\frac{\alpha+\beta\omega}{\gamma+\delta\omega}\right),
\]
and then interchange $\alpha, \beta, \gamma, \delta$ with $-\gamma, -\delta, \alpha, \beta$, there results
\weq{7}{
f_1\!\left(\frac{\gamma+\delta\omega}{\alpha+\beta\omega}\right) = \left(\frac{2}{\gamma}\right)\varrho\, e^{-\frac{3\pi i}{8}\gamma(2\alpha-\delta)}\, f_2(\omega), \qquad \delta \equiv 0 \pmod{2},
}
\weq{8}{
{}= \left(\frac{2}{\delta}\right)\varrho\, e^{-\frac{3\pi i}{8}\delta(2\beta+\gamma)}\, f_1(\omega), \qquad \gamma \equiv 0 \pmod{2},
}
\weq{9}{
{}= -\varrho\, e^{-\frac{3\pi i}{8}\delta(2\beta+\gamma)}\, f(\omega), \qquad \gamma-\delta \equiv 0 \pmod{2}.
}

From (9) and (6) one obtains, on replacing $\omega$ by $\dfrac{-\gamma+\alpha\omega}{\delta-\beta\omega}$ and then $\alpha, \beta, \gamma, \delta$ by $\delta, -\beta, -\gamma, \alpha$:
\weq{10}{
f\!\left(\frac{\gamma+\delta\omega}{\alpha+\beta\omega}\right) = -\varrho\, e^{-\frac{3\pi i}{8}\alpha(2\beta+\gamma)}\, f_1(\omega), \qquad \alpha-\gamma \equiv 0 \pmod{2},
}
\weq{11}{
f\!\left(\frac{\gamma+\delta\omega}{\alpha+\beta\omega}\right) = -\varrho\, e^{\frac{3\pi i}{8}\beta(2\alpha-\delta)}\, f_2(\omega), \qquad \beta-\delta \equiv 0 \pmod{2},
}
and if finally in (10) one replaces $\omega$ by $\omega+9$ and $\gamma, \alpha$ by $\gamma-9\delta, \alpha-9\beta$:
\weq{12}{
f\!\left(\frac{\gamma+\delta\omega}{\alpha+\beta\omega}\right) = \varrho\left(\frac{2}{\alpha-\beta}\right) e^{-\frac{3\pi i}{8}(\alpha-\beta)(\alpha+\beta+\gamma-\delta)}\, f(\omega),
}
\[
\alpha+\beta+\gamma-\delta \equiv 0 \pmod{2}.
\]
\footnote{The functions $\varphi(\omega), \psi(\omega), \chi(\omega)$ introduced by \gesperrt{Hermite} (Sur la théorie des équations modulaires, Paris 1859) are connected with the functions $f(\omega), f_1(\omega), f_2(\omega)$ by the equations
\[
f(\omega) = \frac{\sqrt[6]{2}}{\chi(\omega)}, \quad f_1(\omega) = \sqrt[6]{2}\,\frac{\psi(\omega)}{\chi(\omega)}, \quad f_2(\omega) = \sqrt[6]{2}\,\frac{\varphi(\omega)}{\chi(\omega)}
\]
shown above. The transformation formulae of the $f$-functions can be recast in manifold ways by using the relation $\alpha\delta-\beta\gamma=1$. Thus the formulae given by \gesperrt{Hermite} loc.\,cit.\ are not at once to be recognised as identical with ours. A simple calculation, however, shows their agreement. The formulae given above have the advantage that, without any loss of simplicity, they combine two of the six transformation classes at a time into a single expression.}

\begin{center}\rule{2cm}{0.4pt}\end{center}
